\documentclass{article}
\usepackage[utf8]{inputenc}
\usepackage[T1]{fontenc}
\usepackage[a4paper, top=2.2cm, bottom=2.2cm, left=1.9cm, right=1.9cm, columnsep=0.75cm]{geometry}
\usepackage{multicol}
\usepackage{times}
\usepackage{amsmath}
\usepackage{amssymb}
\usepackage{graphicx}
\usepackage{booktabs}
\usepackage{array}
\usepackage{caption}
\usepackage{hyperref}
\usepackage{cite}
\usepackage{enumitem}
\usepackage{titlesec}
\usepackage{fancyhdr}
\usepackage{xcolor}
\usepackage{microtype}

\titleformat{\section}{\normalsize\bfseries\centering}{\thesection.}{0.5em}{\MakeUppercase}
\titleformat{\subsection}{\normalsize\itshape}{\thesubsection.}{0.5em}{}
\titlespacing{\section}{0pt}{7pt}{3pt}
\titlespacing{\subsection}{0pt}{5pt}{2pt}

\setlength{\parindent}{1em}
\setlength{\parskip}{0pt}

\pagestyle{fancy}
\fancyhf{}
\renewcommand{\headrulewidth}{0pt}
\fancyfoot[C]{\small TER encadr\'e par Arnaud Labourel et Fran\c{c}ois-Xavier Dup\'e.}

\begin{document}

% ─── TITRE ────────────────────────────────────────────────────────────────────
\begin{center}
{\large\bfseries \'Etude de la robustesse des agr\'egateurs en apprentissage f\'ed\'er\'e\\
face aux attaques byzantines avec ByzFL}

\medskip
{\normalsize Arezki C\'elia, Abbes Samy}\\
{\small Parcours IAAA --- 2025--2026}
\end{center}

\medskip

% ─── RÉSUMÉ ───────────────────────────────────────────────────────────────────
\noindent\textbf{R\'esum\'e} ---
L'apprentissage f\'ed\'er\'e permet d'entra\^iner un mod\`ele global \`a partir de donn\'ees distribu\'ees
chez plusieurs clients, sans centralisation des donn\'ees brutes. Ce cadre est cependant vuln\'erable
aux \textit{attaques byzantines non cibl\'ees}, dont l'objectif est de d\'egrader ou de casser compl\`etement
le mod\`ele global. Ce travail compare, \`a l'aide de la biblioth\`eque ByzFL, douze agr\'egateurs robustes
face \`a huit attaques, sur le jeu de donn\'ees MNIST en distributions IID et Non-IID. Deux exp\'eriences
particuli\`erement \'eclairantes sont analys\'ees en d\'etail~: l'attaque IPM en IID et l'attaque Gaussian
en Non-IID. Nos r\'esultats montrent qu'aucun agr\'egateur n'est universellement optimal et que
l'h\'et\'erog\'en\'eit\'e Non-IID amplifie consid\'erablement la difficult\'e du probl\`eme.

\smallskip
\noindent\textbf{Mots cl\'es} --- apprentissage f\'ed\'er\'e, attaques byzantines, agr\'egation robuste,
IID, Non-IID, ByzFL, MNIST

\bigskip

% ─── CORPS À DEUX COLONNES ────────────────────────────────────────────────────
\begin{multicols}{2}

%=======================================================================
\section{Introduction}

L'apprentissage f\'ed\'er\'e (\textit{Federated Learning}, FL) est un paradigme d'apprentissage
automatique distribu\'e propos\'e par McMahan \textit{et al.}~\cite{mcmahan2017}. Son principe
fondateur est simple, plut\^ot que de centraliser les donn\'ees sur un serveur, on laisse
chaque participant, appel\'e \textit{client}, entra\^iner un mod\`ele localement sur ses propres
donn\'ees, puis envoyer uniquement les param\`etres appris (gradients ou poids) \`a un serveur central
charg\'e de les agr\'eger. Le mod\`ele global est ainsi mis \`a jour sans que les donn\'ees brutes
ne quittent jamais les appareils des participants.

Ce paradigme est particuli\`erement int\'eressant dans des contextes o\`u la confidentialit\'e
des donn\'ees est cruciale comme la sant\'e (dossiers m\'edicaux), la t\'el\'ephonie mobile (historique de
navigation, saisie au clavier), ou la finance. Il permet de profiter de la richesse collective
des donn\'ees distribu\'ees sans violer la vie priv\'ee des utilisateurs.

Cependant, cette architecture distribu\'ee introduit une vuln\'erabilit\'e fondamentale. Le
serveur ne peut pas v\'erifier le contenu des mises \`a jour re\c{c}ues. Un client malveillant, dit \textit{byzantin}, peut donc envoyer des gradients arbitrairement corrompus pour
d\'evier ou d\'etruire l'apprentissage. Ce d\'efi s'inscrit dans le cadre g\'en\'eral du \textit{probl\`eme
du consensus} en syst\`emes distribu\'es~\cite{lamport1982}, o\`u le serveur doit obtenir un accord correct et convergent malgr\'e la pr\'esence de participants hostiles. Cette recherche d'un consensus robuste constitue l'objectif centrale \'etudi\'ee dans ce travail.

Ce rapport pr\'esente les exp\'eriences r\'ealis\'ees avec la biblioth\`eque ByzFL~\cite{byzfl2023},
en comparant syst\'ematiquement douze agr\'egateurs robustes face \`a huit types d'attaques,
sur le jeu de donn\'ees MNIST, en distributions IID et Non-IID.
%=======================================================================
\section{Apprentissage F\'ed\'er\'e}

\subsection{Fonctionnement g\'en\'eral}

Dans un syst\`eme d'apprentissage f\'ed\'er\'e classique, l'objectif est d'optimiser un mod\`ele d'apprentissage automatique, comme un r\'eseau de neurones adapt\'e \`a la classification d'images. Le mod\`ele est d\'efini par un vecteur de param\`etres math\'ematiques appel\'es les poids du mod\`ele et not\'es $w$. L'apprentissage consiste \`a minimiser une fonction de co\^ut globale qui mesure l'erreur entre les pr\'edictions du mod\`ele et la r\'ealit\'e.

Le protocole se d\'eroule en plusieurs iterations successives. \`A chaque iteration $t$, le serveur distribue les param\`etres du mod\`ele global courant $w^t$ \`a l'ensemble des $n$ clients. Chaque client $i$ \'evalue l'erreur du mod\`ele sur ses propres donn\'ees locales. Il calcule ensuite un gradient $g_i^t$. Ce gradient correspond \`a la d\'eriv\'ee de la fonction d'erreur. Plus simplement, c'est un vecteur qui indique tr\`es pr\'ecis\'ement dans quelle direction et avec quelle force il faut modifier les poids du mod\`ele pour r\'eduire l'erreur locale. Le client envoie ensuite ce vecteur au serveur.

Celui-ci agr\`ege les contributions re\c{c}ues selon une r\`egle math\'ematique $\mathcal{A}$ pour produire la mise \`a jour globale
\[
  w^{t+1} = w^t - \eta \cdot \mathcal{A}(g_1^t, \ldots, g_n^t)
\]
o\`u $\eta$ est le taux d'apprentissage. Le mod\`ele converge progressivement vers une solution exploitant l'ensemble des donn\'ees distribu\'ees, sans jamais avoir besoin de les centraliser.

\subsection{Vuln\'erabilit\'e aux clients byzantins}

Parmi les $n$ clients, $f$ peuvent \^etre byzantins. Ces participants hostiles sont libres d'envoyer n'importe quel vecteur, ind\'ependamment de leurs donn\'ees r\'eelles ou m\^eme de toute donn\'ee. L'hypoth\`ese classique impose $f < n/2$, assurant ainsi une majorit\'e honnête~\cite{blanchard2017}. Si l'agr\'egateur est une simple moyenne (\textit{Average}), m\^eme un seul client byzantin peut potentiellement envoyer un gradient de norme arbitrairement grande et d\'evier compl\`etement l'apprentissage. C'est cette fragilit\'e qui motive l'\'etude des agr\'egateurs robustes.

\subsection{Deux grandes familles d'attaques}

En apprentissage f\'ed\'er\'e, les attaques byzantines se divisent en deux grandes familles. Les \textbf{attaques cibl\'ees}, \'egalement connues sous le nom d'attaques adversarielles ou \textit{backdoor attacks}, cherchent \`a faire classifier correctement l'ensemble des donn\'ees normales tout en introduisant subtilement un comportement erron\'e sur un sous-ensemble pr\'ecis de donn\'ees poss\'edant un marqueur secret. Les \textbf{attaques non cibl\'ees} visent un objectif plus brutal en cherchant \`a d\'egrader ou casser enti\`erement le mod\`ele global, sans se soucier de maintenir une apparence de normalit\'e.

Dans ce travail, nous nous concentrons exclusivement sur les attaques non cibl\'ees. L'ensemble des huit attaques \'etudi\'ees cherchent \`a emp\^echer la convergence du mod\`ele ou \`a en r\'eduire dramatiquement les performances en corrompant les mises \`a jour envoy\'ees au serveur.

%=======================================================================
\section{Jeu de Donn\'ees et Distributions}

\subsection{MNIST}

Le jeu de donn\'ees MNIST (\textit{Modified National Institute of Standards and
Technology})~\cite{lecun1998} est le support exp\'erimental de toutes les exp\'eriences
men\'ees dans ce travail. Il se compose de \textbf{70~000 images en niveaux de gris}
repr\'esentant des chiffres manuscrits de 0 \`a 9, soit \textbf{10 classes}. Chaque
image est de taille \textbf{28$\times$28 pixels}, repr\'esent\'ee comme un vecteur
de 784 valeurs entre 0 et 255.

La r\'epartition est la suivante~: \textbf{60~000 images d'entra\^inement} et \textbf{10~000
images de test}. Les classes sont quasiment \'equilibr\'ees~, environ \textbf{6~000 images
par classe} dans l'ensemble d'entra\^inement, et environ \textbf{1~000 images par classe}
dans l'ensemble de test.

MNIST constitue un choix id\'eal pour ce type d'\'etude~, il est suffisamment simple
pour permettre des entra\^inements rapides sur de nombreuses configurations, bien connu
et r\'ef\'erenc\'e dans la litt\'erature, et suffisamment riche pour mettre en \'evidence
des diff\'erences significatives entre strat\'egies de d\'efense. Sa simplicit\'e garantit
que les \'ecarts de performances observ\'es entre agr\'egateurs refl\`etent bien l'effet
des attaques, et non une difficult\'e intrins\`eque de la t\^ache de classification.

Un jeu de donn\'ees plus complexe, comme \textbf{CIFAR-10}~\cite{krizhevsky2009},
compos\'e de 60~000 images couleur RGB 32$\times$32 en 10 classes, serait davantage
adapt\'e \`a un r\'eseau de neurones convolutionnel (CNN). Les CNN exploitent la structure
spatiale des images gr\^ace \`a des couches de convolution, de pooling et d'activation,
et obtiennent des performances nettement sup\'erieures sur des t\^aches visuelles complexes.
Dans une suite \`a ce travail, CIFAR-10 associ\'e \`a un CNN constituerait une \'etape naturelle
pour tester la robustesse des agr\'egateurs dans un contexte complexifié.

\subsection{Distributions IID et Non-IID}

La mani\`ere dont les donn\'ees sont r\'eparties entre les clients a une influence
d\'eterminante sur la difficult\'e du probl\`eme.

En configuration \textbf{IID} (\textit{Independent and Identically Distributed}),
chaque client re\c{c}oit un \'echantillon al\'eatoire et uniform\'ement distribu\'e de
l'ensemble des donn\'ees. Chaque classe est ainsi repr\'esent\'ee de mani\`ere \'equilibr\'ee
chez chaque client. Les gradients honnêtes produits sont donc similaires entre clients
et pointent naturellement dans la m\^eme direction. Dans ce contexte, d\'etecter une
mise \`a jour byzantine est relativement accessible~, elle se d\'emarque du consensus.

En configuration \textbf{Non-IID}, la r\'epartition est biais\'ee~, chaque client
ne poss\`ede que certaines classes, ou les observe avec des fr\'equences tr\`es in\'egales.
Par exemple, un client peut n'avoir re\c{c}u que des images du chiffre 3 et du chiffre 7,
tandis qu'un autre client aura exclusivement des 0 et des 1. Ce biais de distribution
cr\'ee une h\'et\'erog\'en\'eit\'e profonde entre clients~, leurs gradients honnêtes
divergent naturellement, m\^eme sans aucune attaque. Cette dispersion l\'egitime
complique drastiquement le travail des agr\'egateurs robustes, qui doivent contourner les perturbations introduites par les
byzantins en faisant face à une variabilit\'e naturelle due au Non-IID 

%=======================================================================
\section{Attaques et Agr\'egateurs}

\subsection{Attaques byzantines \'etudi\'ees}

Les huit attaques non cibl\'ees \'etudi\'ees visent toutes \`a d\'egrader le mod\`ele global,
mais par des m\'ecanismes tr\`es diff\'erents.

\textbf{Sign Flipping (SF)}~\cite{li2019}~: l'attaquant inverse simplement le signe de son gradient
r\'eel~: $\tilde{g} = -g_i$. L'effet est brutal~: le mod\`ele est pouss\'e syst\'ematiquement
dans la mauvaise direction \`a chaque iteration. Simple mais redoutablement efficace contre
les agr\'egateurs qui ne filtrent pas les valeurs extr\^emes.

\textbf{Inner Product Manipulation (IPM)}~\cite{xie2020}~: plut\^ot que d'agir brutalement,
l'attaquant envoie un gradient dont le produit scalaire avec la vraie moyenne des gradients
honnêtes est fortement n\'egatif. Cela d\'evie l'optimisation de mani\`ere cibl\'ee et structur\'ee,
rendant l'attaque plus difficile \`a d\'etecter.

\textbf{Optimal IPM (Opt-IPM)}~: version optimis\'ee d'IPM, qui maximise analytiquement
l'angle entre le gradient malveillant et la moyenne honnête, tout en restant discret.

\textbf{A Little Is Enough (ALIE)}~\cite{baruch2019}~: attaque particuli\`erement
sournoise. L'attaquant analyse la distribution statistique des gradients honnêtes et
injecte une perturbation calibr\'ee pour se fondre dans leur bruit naturel. L'attaque
est "invisible" aux yeux de nombreuses d\'efenses~, le gradient malveillant est
statistiquement plausible, mais oriente subtilement l'apprentissage dans la mauvaise direction.

\textbf{Optimal ALIE (Opt-ALIE)}~: variante d'ALIE qui optimise le vecteur d'attaque
pour maximiser son impact tout en \'evitant d'\^etre \'elimin\'e par les agr\'egateurs
bas\'es sur la norme ou la distance.

\textbf{Mimic}~: l'attaquant copie exactement le gradient d'un client honnête. En se
d\'eguisant en participant l\'egitime, il trompe les d\'efenses bas\'ees sur la similarit\'e qui ne peuvent plus le diff\'erencier d'un client normal. Cette attaque
r\'ev\`ele une faiblesse profonde des m\'ethodes de s\'election par proximit\'e.

\textbf{Gaussian}~: le gradient malveillant est tir\'e d'une distribution gaussienne
centr\'ee sur z\'ero, independante de toute donn\'ee r\'eelle. L'attaque est brute et
non structur\'ee~, elle injecte un bruit al\'eatoire. Si elle semble na\"{i}ve, elle se
r\'ev\`ele d\'evastatrice pour les agr\'egateurs sans m\'ecanisme de filtrage, et ses
effets sont amplifi\'es par le contexte Non-IID.

\textbf{Infinity}~: les gradients malveillants ont une norme extr\^emement \'elev\'ee,
visant \`a dominer la somme agr\'eg\'ee et \`a \'ecraser les contributions honnêtes.

\subsection{Agr\'egateurs compar\'es}

Douze agr\'egateurs sont \'evalu\'es, des plus simples aux plus sophistiqu\'es.

\textbf{Average}~: moyenne arithm\'etique na\"{i}ve de tous les gradients. Aucune
r\'esistance aux attaques~. Un seul byzantin suffit \`a d\'evier l'ensemble de l'apprentissage.
Sert de r\'ef\'erence basse.

\textbf{Median}~\cite{yin2018}~: m\'ediane coordonn\'ee par coordonn\'ee. Plus robuste
qu'Average, car les valeurs extr\^emes ont moins d'influence, mais vuln\'erable aux
attaques qui ciblent la m\'ediane elle-m\^eme.

\textbf{TrMean}~\cite{yin2018}~: moyenne tronqu\'ee. \'Elimine les $f$ gradients les
plus grands et les plus petits sur chaque coordonn\'ee avant de faire la moyenne. Sa robustesse dépend du nombre de valeurs extrêmes supprimées avant le calcul de la moyenne.

\textbf{Geometric Median}~\cite{chen2017}~: point de l'espace qui minimise la somme
des distances euclidiennes \`a l'ensemble des gradients. Tr\`es robuste aux vecteurs
\'eloign\'es de la majorit\'e, mais co\^uteux \`a calculer.

\textbf{Krum}~\cite{blanchard2017}~: s\'electionne le gradient dont la somme des
distances aux $n-f-2$ plus proches voisins est minimale. Il identifie ainsi le
gradient le plus repr\'esentatif de la majorit\'e. Tr\`es efficace contre les attaques
directes, mais mis en d\'efaut par Mimic.

\textbf{Multi-Krum}~: variante de Krum s\'electionnant les $m > 1$ meilleurs gradients
au sens du score Krum, puis en faisant leur moyenne. Offre un meilleur compromis
biais-variance que Krum seul.

\textbf{Centered Clipping (CC)}~\cite{karimireddy2021}~: projette chaque gradient dans
une boule de rayon $\tau$ centr\'ee autour d'une estimation du centre honnête. Tr\`es
efficace en Non-IID, car il adapte son centre aux donn\'ees h\'et\'erog\`enes.

\textbf{MDA} (\textit{Minimum Distance Aggregation})~: s\'electionne le sous-ensemble
de gradients dont la dispersion interne est minimale, en postulant que les honnêtes
sont plus concentr\'es que les byzantins.

\textbf{MoNNA}~: agr\'egateur bas\'e sur le momentum~, il utilise l'historique des mises
\`a jour pour d\'etecter les gradients inconsistants avec la dynamique observ\'ee.

\textbf{Meamed}~: variante de la m\'ediane qui utilise une estimation de la moyenne pour donner plus ou moins d'importance aux gradients, afin de combiner robustesse et efficacit\'e.

\textbf{CAF}~: filtre les gradients par blocs de coordonn\'ees, cherchant \`a isoler
les dimensions les plus corrompues.

\textbf{SMEA}~: m\'ethode r\'ecente qui s\'electionne les gradients jug\'es fiables avant de calculer leur moyenne, afin de limiter l'influence des clients byzantins, montrant
de tr\`es bonnes performances sur les attaques gaussiennes et Mimic en Non-IID.

%=======================================================================
\section{Protocole Exp\'erimental}

Les exp\'eriences ont \'et\'e r\'ealis\'ees avec \textbf{5 clients honnêtes} et
\textbf{2 clients byzantins}, sur le jeu de donn\'ees MNIST, avec plusieurs iterations
d'entra\^inement (jusqu’à 500 itérations pour les expériences longues sur Google Colab avec GPU, que nous avons réduites à environ 200 vers la fin par manque de temps et de ressources). Chaque sc\'enario est d\'efini dans un fichier \textbf{JSON} sp\'ecifiant
l'attaque, l'agr\'egateur, la distribution (IID ou Non-IID) et le nombre d'iterations.
Un moteur Python automatise le lancement et sauvegarde les r\'esultats en CSV pour
visualisation. Certains sc\'enarios ont \'et\'e r\'ep\'et\'es sur plusieurs seeds al\'eatoires (5 runs)
pour stabiliser les r\'esultats.

%=======================================================================
%=======================================================================
\section{R\'esultats Exp\'erimentaux}

Parmi l'ensemble des exp\'eriences men\'ees, nous avons cherch\'e \`a identifier la meilleure attaque et le meilleur agr\'egateur selon un crit\`ere de pire cas. Une attaque est consid\'er\'ee comme la meilleure si elle arrive \`a minimiser la pr\'ecision maximale de l'ensemble des agr\'egateurs. \`A l'inverse, le meilleur agr\'egateur est celui qui maximise la pr\'ecision minimale face \`a toutes les attaques confondues.

\subsection{La meilleure attaque et le meilleur agr\'egateur}

L'analyse compl\`ete de toutes les courbes d'apprentissage confirme que l'attaque Optimal Inner Product Manipulation en distribution Non-IID s'impose d\'efinitivement comme la plus redoutable. Elle parvient \`a faire chuter drastiquement les performances d'agr\'egateurs pourtant tr\`es robustes en temps normal. Des m\'ethodes comme SMEA ou Multi-Krum peinent \`a d\'epasser les 50 pour cent de pr\'ecision sous cette attaque. Les exp\'eriences suppl\'ementaires montrent que des attaques alternatives comme Gaussian ou A Little Is Enough sont finalement bien g\'er\'ees par la majorit\'e des d\'efenses. L'attaque Gaussian injecte un bruit trop \'evident qui se retrouve tr\`es vite bloqu\'e par les algorithmes. L'attaque A Little Is Enough est plus sournoise mais elle ne r\'eussit pas \`a d\'etruire le mod\`ele d'apprentissage violement. En effet, il n'y a que l'aggrégateur Meamed qui peine à contrer cette attaque en distribution non-iid. Le biais naturel des donn\'ees Non-IID reste le facteur aggravant qui complique la t\^ache des d\'efenses et permet aux meilleures attaques de faire un maximum de d\'eg\^ats.

Du c\^ot\'e des d\'efenses, la moyenne tronqu\'ee (\textit{TrMean}) et la m\'ediane g\'eom\'etrique (\textit{Geometric Median}) s'en sortent le mieux de fa\c{c}on globale. M\^eme lorsque les exp\'eriences sont prolong\'ees jusqu'\`a 500 it\'erations pour tester la r\'esistance sur le long terme, ces m\'ethodes prouvent leur grande fiabilit\'e. Elles ne finissent pas forc\'ement premi\`eres sur chaque graphique individuel, mais elles ne s'effondrent jamais. Elles assurent ainsi une performance minimale tr\`es robuste sur toutes les exp\'eriences confondues, l\`a o\`u d'autres agr\'egateurs tombent parfois \`a des scores proches du hasard sur des attaques sp\'ecifiques. Il convient cependant de noter que la s\'ev\'erit\'e de l'attaque \textit{Optimal IPM} parvient \`a affaiblir significativement les capacit\'es de ces deux d\'efenses r\'esilientes, alors que des agr\'egateurs plus sophistiqu\'es, tels que \textit{MoNNA}, \textit{MDA} et \textit{CAF}, ont su d\'emontrer leur efficacit\'e dans ce contexte.

\subsection{Exp\'erience 1 avec IPM en distribution IID}

L'attaque Inner Product Manipulation (IPM) reste assez simple dans son approche algorithmique. L'attaquant construit un gradient dont le produit scalaire avec la direction honn\^ete est n\'egatif pour d\'evier le mod\`ele. L'effet est subtil mais profond. Plut\^ot que de d\'etruire brutalement les mises \`a jour, l'attaque d\'evie progressivement le mod\`ele dans une mauvaise direction.

Dans cette configuration IID, l'agr\'egateur Multi-Krum atteint la meilleure pr\'ecision finale avec 0.9672. Ce r\'esultat est remarquable, mais il est particuli\`erement int\'eressant de noter que l'agr\'egateur na\"if Average r\'esiste \'etonnamment bien \`a cette attaque sp\'ecifique en IID, parvenant \`a maintenir pr\`es de 95\% de pr\'ecision (0.9491). L'agr\'egateur Centered Clipping offre d'ailleurs des performances similaires (0.9491). Le gain de Multi-Krum reste n\'eanmoins significatif car il r\'eussit \`a optimiser l'entra\^inement au maximum, \'ecartant avec plus de pr\'ecision les vecteurs byzantins pour fr\^oler les 97\%.

Ce succ\`es s'explique par la r\'epartition des donn\'ees. En IID, les gradients honn\^etes sont tr\`es similaires entre eux. Multi-Krum calcule les distances entre les vecteurs et identifie naturellement le groupe majoritaire. Il \'ecarte facilement les byzantins car ils s'\'eloignent trop du consensus.

\subsection{Exp\'erience 2 avec Gaussian en distribution Non-IID}

L'attaque Gaussian injecte simplement du bruit al\'eatoire. C'est l'attaque la plus na\"ive possible.

En distribution Non-IID, l'agr\'egateur SMEA obtient la meilleure performance avec une pr\'ecision de 0.9695. En distribution IID sur la m\^eme attaque Gaussian, SMEA reste d'ailleurs en t\^ete (0.9717), talonn\'e de tr\`es pr\`es par Multi-Krum (0.9715) et Geometric Median (0.971). Le contraste le plus frappant concerne la moyenne classique (Average), alors qu'elle parvenait \`a maintenir environ 95 pour cent face \`a IPM, elle s'effondre compl\`etement sous l'attaque Gaussian, tombant \`a 22.3\% en IID et \`a environ 10\% en Non-IID (soit l'\'equivalent d'une classification totalement au hasard).

Le point le plus int\'eressant est de voir \`a quel point la robustesse d\'epend du contexte. Par exemple, Multi-Krum brille face aux attaques en IID, mais son efficacit\'e chute drastiquement face \`a IPM en Non-IID (tombant \`a moins de 59 \% de pr\'ecision). La distribution Non-IID amplifie dramatiquement la difficult\'e de l'apprentissage. Les gradients honn\^etes sont d\'ej\`a tr\`es dispers\'es \`a cause de la r\'epartition in\'egale des donn\'ees. Les agr\'egateurs ont alors \'enorm\'ement de mal \`a s\'eparer la variabilit\'e l\'egitime des perturbations malveillantes.

\begin{center}
\captionof{table}{Accuracy finale compar\'ee sur IPM et Gaussian.}
\label{tab:results}
\footnotesize
\setlength{\tabcolsep}{3pt} % Réduit pour faire tenir 5 colonnes
\begin{tabular}{lcccc}
\toprule
\textbf{Agr\'egateur} & \textbf{IPM} & \textbf{IPM} & \textbf{Gaussian} & \textbf{Gaussian} \\
 & \textbf{(IID)} & \textbf{(Non-IID)} & \textbf{(IID)} & \textbf{(Non-IID)} \\
\midrule
Average          & 0.94912       & 0.9473        & 0.223         & $\sim$0.10  \\
Geometric Median & 0.9636        & 0.9375        & 0.971         & 0.967     \\
Multi-Krum       & \textbf{0.9672} & 0.5895      & 0.9715        & 0.9644         \\
CC               & 0.9491        & 0.9472        & 0.9518        & 0.94008        \\
SMEA             & 0.9501        & 0.4609        & \textbf{0.9717} & \textbf{0.9695} \\
\bottomrule
\end{tabular}
\end{center}
\vspace{0.2cm}

%=======================================================================
\subsection{Absence d'agr\'egateur universel}

Le constat le plus saillant de ce travail est qu'aucun agr\'egateur ne domine toutes les situations. Multi Krum excelle par exemple face \`a l'attaque IPM en distribution IID mais reste tr\`es faible dans un contexte Non IID. Les deux graphiques plac\'es juste en dessous permettent de comparer visuellement cette chute de performance lors du passage d'une distribution \`a l'autre.

\begin{center}
\begin{minipage}{0.5\linewidth}
\centering
\includegraphics[width=\linewidth]{iid_innerproduct_multikrum.png}
\captionof{figure}{Multi Krum avec IPM(IID)}
\label{fig_multi_krum_iid}
\end{minipage}%
\hfill%
\begin{minipage}{0.5\linewidth}
\centering
\includegraphics[width=\linewidth]{non_iid_innerproduct_multikrum.png}
\captionof{figure}{Multi Krum avec IPM(Non IID)}
\label{fig_multi_krum_non_iid}
\end{minipage}
\end{center}

SMEA est optimal face \`a Gaussian en Non-IID mais se r\'ev\`ele tr\`es fragile sur d'autres attaques tel que \textit{IPM}. Le choix d'un agr\'egateur doit imperativement d\'ependre d'une analyse du profil de menace pr\'evu dans le contexte de d\'eploiement.

\subsection{Impact du cadre Non-IID}

La distribution Non-IID repr\'esente syst\'ematiquement le sc\'enario le plus exigeant. En IID, les gradients honn\^etes forment un groupe compact ce qui rend la d\'etection des byzantins plus \'evidente. En Non-IID, ce groupe s'\'etend. Les agr\'egateurs bas\'es sur la concentration comme Krum ou Multi-Krum \'eliminent parfois des gradients honn\^etes l\'egitimement atypiques, ce qui d\'egrade leurs propres capacit\'es d'apprentissage.

\subsection{Limites du travail}

Toutes les attaques n'ont pas pu \^etre \'evalu\'ees avec la m\^eme profondeur. En particulier, l'attaque \textit{Optimal ALIE} a d\^u \^etre \'ecart\'ee des exp\'eriences longues, son temps d'ex\'ecution d\'emesur\'e provoquant des \'echecs syst\'ematiques sur Google Colab. Cette lourdeur s'explique par la m\'ecanique de l'attaque, qui \`a chaque it\'eration, doit r\'esoudre un co\^uteux probl\`eme d'optimisation sous contraintes pour calibrer parfaitement sa perturbation sur la distribution des gradients honn\^etes. Par ailleurs, nos r\'esultats reposent sur MNIST, un jeu de donn\'ees classique mais simple. Des exp\'eriences futures sur CIFAR-10 avec un r\'eseau convolutif pourraient affiner ces conclusions. Enfin, les param\`etres des agr\'egateurs ont \'et\'e utilis\'es avec leurs valeurs par d\'efaut, sans r\'eglage sp\'ecifique par attaque.

%=======================================================================
\section{Conclusion}

Ce TER a permis de mettre en place un pipeline exp\'erimental automatis\'e pour comparer douze agr\'egateurs robustes face \`a huit attaques byzantines non cibl\'ees. Ces simulations ont \'et\'e r\'ealis\'ees sur MNIST avec des distributions IID et Non-IID.

L'int\'er\^et de l'apprentissage f\'ed\'er\'e r\'eside dans sa capacit\'e \`a exploiter des donn\'ees distribu\'ees sans les centraliser. Sa vuln\'erabilit\'e aux clients byzantins est bien r\'eelle. La moyenne classique tombe \`a 10 pour cent de pr\'ecision sous de nombreuses attaques, l\`a o\`u des d\'efenses adapt\'ees maintiennent d'excellents scores. Le choix de l'agr\'egateur est d\'eterminant et aucune m\'ethode n'est universellement parfaite. L'h\'et\'erog\'en\'eit\'e des donn\'ees en Non-IID rend toujours le probl\`eme bien plus difficile \`a r\'esoudre.

La biblioth\`eque ByzFL constitue un outil solide pour poursuivre ces investigations. Il sera possible d'ajouter de nouvelles configurations d'attaques ou d'utiliser des mod\`eles d'intelligence artificielle plus complexes pour prolonger ce travail de recherche.
%=======================================================================
\section{Bilan Personnel}

\subsection*{Bilan de Samy}

Ce projet a \'et\'e pour moi une d\'ecouverte concr\`ete de ce que signifie mener
une \'etude exp\'erimentale scientifique. J'avais des notions th\'eoriques sur
l'apprentissage automatique, mais l'exp\'erience de construire un pipeline complet, de la d\'efinition des sc\'enarios JSON jusqu'\`a la g\'en\'eration des courbes
d'analyse, m'a apport\'e une vision bien plus globale du processus de recherche.
J'ai particuli\`erement appr\'eci\'e le travail d'analyse des r\'esultats, qui m'a oblig\'e \`a r\'efl\'echir au-del\`a des
chiffres et \`a formuler des explications formelle. Ce TER m'a aussi enseign\'e
la valeur de l'organisation~: structurer un d\'ep\^ot, nommer proprement les fichiers
de r\'esultats, r\'ep\'eter les exp\'eriences sur plusieurs seeds, autant d'habitudes
que je consid\`ere d\'esormais indispensables dans tout travail scientifique.

\subsection*{Bilan de C\'elia}

Ce TER a \'et\'e l'occasion pour moi d'aborder pour la premi\`ere fois la probl\'ematique
de la s\'ecurit\'e en apprentissage automatique, et plus pr\'ecis\'ement en apprentissage
f\'ed\'er\'e. Avant ce projet, je n'avais qu'une connaissance superficielle de ce domaine.
Travailler sur les attaques byzantines et les m\'ecanismes de d\'efense m'a permis
de comprendre que la robustesse est un probl\`eme fondamentalement contextuel, autrement dit, il n'existe pas de solution miracle, et chaque agr\'egateur a ses forces et ses failles
selon le sc\'enario. J'ai \'egalement d\'evelopp\'e des comp\'etences pratiques en
automatisation d'exp\'eriences et en analyse de courbes. La difficult\'e principale
a \'et\'e de g\'erer la quantit\'e de configurations \`a tester et d'en tirer des
conclusions claires et justifi\'ees. Ce travail m'a donn\'e envie d'aller plus loin
sur les questions de confidentialit\'e et de s\'ecurit\'e des syst\`emes d'IA distribu\'es.

%=======================================================================
\end{multicols}

\newpage

\begin{center}
{\normalsize\bfseries 9. \MakeUppercase{R\'ef\'erences}}
\end{center}

% On masque le titre généré automatiquement par thebibliography 
% pour conserver ton propre titre personnalisé juste au-dessus.
\renewcommand{\refname}{\vspace{-5ex}}

\begin{multicols}{2}
\small
\begin{thebibliography}{12}
\setlength{\itemsep}{3pt}

  \bibitem{mcmahan2017} H.~B. McMahan, E.~Moore, D.~Ramage, S.~Hampson,
        et B.~A. y~Arcas, ``Communication-efficient learning of deep networks
        from decentralized data,'' in \textit{Proc. AISTATS}, 2017.

  \bibitem{lamport1982} L.~Lamport, R.~Shostak, et M.~Pease, ``The Byzantine
        generals problem,'' \textit{ACM Trans. Program. Lang. Syst.}, vol.~4,
        no.~3, pp.~382--401, 1982.

  \bibitem{blanchard2017} P.~Blanchard, E.~M.~El~Mhamdi, R.~Guerraoui,
        et J.~Stainer, ``Machine learning with adversaries: Byzantine tolerant
        gradient descent,'' in \textit{Proc. NeurIPS}, 2017.

  \bibitem{chen2017} Y.~Chen, L.~Su, et J.~Xu, ``Distributed statistical
        machine learning in adversarial settings: Byzantine gradient descent,''
        \textit{Proc. ACM Meas. Anal. Comput. Syst.}, vol.~1, no.~2, 2017.

  \bibitem{karimireddy2021} S.~P. Karimireddy, L.~He, et M.~Jaggi,
        ``Byzantine-robust learning on heterogeneous datasets via bucketing,''
        in \textit{Proc. ICLR}, 2021.

  \bibitem{li2019} L.~Li, W.~Xu, T.~Chen, G.~B.~Giannakis, et Q.~Ling,
        ``RSA: Byzantine-robust stochastic aggregation methods for distributed
        learning from heterogeneous datasets,'' in \textit{Proc. AAAI}, 2019.

  \bibitem{xie2020} C.~Xie, O.~Koyejo, et I.~Gupta, ``Fall of empires:
        Breaking byzantine-tolerant SGD by inner product manipulation,''
        in \textit{Proc. UAI}, 2020.

  \bibitem{baruch2019} G.~Baruch, M.~Baruch, et Y.~Goldberg,
        ``A little is enough: Circumventing defenses for distributed learning,''
        in \textit{Proc. NeurIPS}, 2019.

  \bibitem{yin2018} D.~Yin, Y.~Chen, R.~Kannan, et P.~Bartlett,
        ``Byzantine-robust distributed learning: Towards optimal statistical
        rates,'' in \textit{Proc. ICML}, 2018.

  \bibitem{lecun1998} Y.~LeCun, L.~Bottou, Y.~Bengio, et P.~Haffner,
        ``Gradient-based learning applied to document recognition,''
        \textit{Proc. IEEE}, vol.~86, no.~11, pp.~2278--2324, 1998.

  \bibitem{krizhevsky2009} A.~Krizhevsky, ``Learning multiple layers of
        features from tiny images,'' Tech. Rep., Univ. Toronto, 2009.

  \bibitem{byzfl2023} El~Mhamdi \textit{et al.}, ``ByzFL: A Python library
        for Byzantine-robust federated learning simulation,''
        \url{https://github.com/LPD-EPFL/ByzFL}, 2023.

\end{thebibliography}
\end{multicols}

\end{document}